% ---------------------------------------------------------------------
\section{Kết quả mô phỏng}
\label{sec:ketqua}
% ---------------------------------------------------------------------

\subsection{Thiết lập mô phỏng}

Toàn bộ hệ thống nằm trong một file LTspice duy nhất \texttt{PIR\_LHI778.asc}: tín hiệu vào, mô hình \LHI{}, nguồn nhiễu, mạch khuếch đại -- lọc và nguồn tham chiếu (hình~\ref{fig:ltspice}). Mô phỏng dùng LTspice 26.1.1 với:
\begin{itemize}
  \item phân tích quá độ \texttt{.tran 0 40 0 0.1m} (40~s, bước tối đa 0,1~ms, đủ để biểu diễn nhiễu 1~kHz; chạy khoảng 15 -- 30~s);
  \item phân tích tần số \texttt{.ac dec 100 0.01 1k} với nguồn thử 1~µW trên phần tử A;
  \item \texttt{.options gmin=1e-15}, \texttt{.save} các nút chính, và các lệnh \texttt{.meas} tự động đo mức nghỉ, biên độ từng kịch bản, nhiễu trước/sau lọc.
\end{itemize}

\begin{landscape}
\begin{figure}[p]
\centering
\begin{subfigure}{\linewidth}\centering
\begin{tikzpicture}[x=1cm,y=1cm,line cap=round,line join=round,wire/.style={line width=0.5pt,c1!75!black},sym/.style={line width=0.45pt,bkred!75!black},lab/.style={font=\fontsize{4.6}{5}\selectfont\sffamily,inner sep=0.5pt},net/.style={font=\fontsize{4.8}{5}\selectfont\sffamily\bfseries,text=c2!85!black,inner sep=0.6pt}]
\draw[wire] (1.447,-2.059) -- (1.447,-2.207);
\draw[wire] (1.447,-1.317) -- (1.447,-1.020);
\draw[wire] (3.524,-2.207) -- (3.524,-1.020);
\draw[wire] (5.602,-1.317) -- (7.679,-1.317);
\draw[wire] (7.679,-1.317) -- (8.421,-1.317);
\draw[wire] (8.421,-1.317) -- (9.015,-1.317);
\draw[wire] (7.679,-1.910) -- (7.679,-2.059);
\draw[wire] (5.602,-3.246) -- (7.679,-3.246);
\draw[wire] (7.679,-3.246) -- (8.421,-3.246);
\draw[wire] (8.421,-3.246) -- (9.015,-3.246);
\draw[wire] (7.679,-3.840) -- (7.679,-3.988);
\draw[wire] (10.350,-1.614) -- (13.466,-1.614);
\draw[wire] (13.466,-1.614) -- (14.208,-1.614);
\draw[wire] (14.208,-1.614) -- (15.098,-1.614);
\draw[wire] (15.098,-1.614) -- (15.989,-1.614);
\draw[wire] (13.466,-2.207) -- (13.466,-2.356);
\draw[wire] (10.350,-2.356) -- (13.466,-2.356);
\draw[wire] (13.466,-2.356) -- (14.208,-2.356);
\draw[wire] (14.208,-2.356) -- (14.653,-2.356);
\draw[wire] (13.466,-2.949) -- (13.466,-3.098);
\draw[wire] (15.098,-2.356) -- (15.098,-3.098);
\draw[wire] (17.324,-2.504) -- (17.324,-3.098);
\draw[wire] (10.350,-3.098) -- (13.466,-3.098);
\draw[wire] (13.466,-3.098) -- (14.208,-3.098);
\draw[wire] (14.208,-3.098) -- (15.098,-3.098);
\draw[wire] (15.098,-3.098) -- (17.324,-3.098);
\draw[wire] (17.324,-3.098) -- (19.550,-3.098);
\draw[wire] (16.434,-1.020) -- (19.550,-1.020);
\draw[wire] (16.434,-1.910) -- (17.324,-1.910);
\draw[wire] (17.324,-1.910) -- (19.550,-1.910);
\draw[wire] (19.550,-1.910) -- (20.144,-1.910);
\draw[wire] (20.144,-1.910) -- (20.885,-1.910);
\draw[wire] (21.627,-1.910) -- (22.369,-1.910);
\fill[c2!85!black] (1.447,-1.020) circle (0.6pt);
\node[net,anchor=south west] at (1.447,-1.020) {PA};
\draw[wire] (1.447,-2.949) -- ++(0,-0.074) ++(-0.148,0) -- ++(0.297,0) -- ++(-0.148,-0.148) -- cycle;
\fill[c2!85!black] (3.524,-1.020) circle (0.6pt);
\node[net,anchor=south west] at (3.524,-1.020) {PB};
\draw[wire] (3.524,-2.949) -- ++(0,-0.074) ++(-0.148,0) -- ++(0.297,0) -- ++(-0.148,-0.148) -- cycle;
\fill[c2!85!black] (9.015,-1.317) circle (0.6pt);
\node[net,anchor=south west] at (9.015,-1.317) {TA};
\draw[wire] (5.602,-2.059) -- ++(0,-0.074) ++(-0.148,0) -- ++(0.297,0) -- ++(-0.148,-0.148) -- cycle;
\draw[wire] (7.679,-2.059) -- ++(0,-0.074) ++(-0.148,0) -- ++(0.297,0) -- ++(-0.148,-0.148) -- cycle;
\draw[wire] (8.421,-2.059) -- ++(0,-0.074) ++(-0.148,0) -- ++(0.297,0) -- ++(-0.148,-0.148) -- cycle;
\fill[c2!85!black] (9.015,-3.246) circle (0.6pt);
\node[net,anchor=south west] at (9.015,-3.246) {TB};
\draw[wire] (5.602,-3.988) -- ++(0,-0.074) ++(-0.148,0) -- ++(0.297,0) -- ++(-0.148,-0.148) -- cycle;
\draw[wire] (7.679,-3.988) -- ++(0,-0.074) ++(-0.148,0) -- ++(0.297,0) -- ++(-0.148,-0.148) -- cycle;
\draw[wire] (8.421,-3.988) -- ++(0,-0.074) ++(-0.148,0) -- ++(0.297,0) -- ++(-0.148,-0.148) -- cycle;
\fill[c2!85!black] (13.466,-1.614) circle (0.6pt);
\node[net,anchor=south west] at (13.466,-1.614) {G};
\fill[c2!85!black] (14.653,-2.356) circle (0.6pt);
\node[net,anchor=south west] at (14.653,-2.356) {M};
\fill[c2!85!black] (19.550,-1.020) circle (0.6pt);
\node[net,anchor=south west] at (19.550,-1.020) {VDD};
\draw[wire] (19.550,-3.098) -- ++(0,-0.074) ++(-0.148,0) -- ++(0.297,0) -- ++(-0.148,-0.148) -- cycle;
\draw[wire] (20.144,-2.652) -- ++(0,-0.074) ++(-0.148,0) -- ++(0.297,0) -- ++(-0.148,-0.148) -- cycle;
\fill[c2!85!black] (20.144,-1.910) circle (0.6pt);
\node[net,anchor=south west] at (20.144,-1.910) {S};
\fill[c2!85!black] (22.369,-1.910) circle (0.6pt);
\node[net,anchor=south west] at (22.369,-1.910) {SN};
\draw[cgray,dashed,line width=0.5pt] (9.608,-0.130) rectangle (19.253,-5.175);
\draw[sym] (1.373,-1.502) -- (1.521,-1.502);
\draw[sym] (1.373,-1.873) -- (1.521,-1.873);
\draw[sym] (1.447,-1.428) -- (1.447,-1.577);
\draw[sym] (1.447,-2.059) -- (1.447,-1.985);
\draw[sym] (1.447,-1.317) -- (1.447,-1.391);
\draw[sym] (1.447,-1.688) ellipse (0.297 and 0.297);
\node[lab,anchor=west] at (1.669,-1.317) {VAC};
\node[lab,anchor=west] at (1.669,-2.059) {AC 1};
\draw[sym] (1.373,-2.393) -- (1.521,-2.393);
\draw[sym] (1.373,-2.764) -- (1.521,-2.764);
\draw[sym] (1.447,-2.319) -- (1.447,-2.467);
\draw[sym] (1.447,-2.949) -- (1.447,-2.875);
\draw[sym] (1.447,-2.207) -- (1.447,-2.281);
\draw[sym] (1.447,-2.578) ellipse (0.297 and 0.297);
\node[lab,anchor=west] at (1.669,-2.207) {BIRA};
\node[lab,anchor=west] at (1.669,-2.949) {V=IRA(time)};
\draw[sym] (3.450,-2.393) -- (3.598,-2.393);
\draw[sym] (3.450,-2.764) -- (3.598,-2.764);
\draw[sym] (3.524,-2.319) -- (3.524,-2.467);
\draw[sym] (3.524,-2.949) -- (3.524,-2.875);
\draw[sym] (3.524,-2.207) -- (3.524,-2.281);
\draw[sym] (3.524,-2.578) ellipse (0.297 and 0.297);
\node[lab,anchor=west] at (3.747,-2.207) {BIRB};
\node[lab,anchor=west] at (3.747,-2.949) {V=IRB(time)};
\draw[sym] (5.602,-1.540) -- (5.639,-1.651);
\draw[sym] (5.602,-1.540) -- (5.565,-1.651);
\draw[sym] (5.565,-1.651) -- (5.639,-1.651);
\draw[sym] (5.602,-1.651) -- (5.602,-1.836);
\draw[sym] (5.602,-2.059) -- (5.602,-1.985);
\draw[sym] (5.602,-1.317) -- (5.602,-1.391);
\draw[sym] (5.602,-1.688) ellipse (0.297 and 0.297);
\node[lab,anchor=west] at (5.824,-1.540) {BQA};
\node[lab,anchor=west] at (5.824,-1.836) {I=ETA*1e-6*V(PA)};
\draw[sym] (7.679,-1.651) -- (7.679,-1.910);
\draw[sym] (7.679,-1.577) -- (7.679,-1.317);
\draw[sym] (7.531,-1.577) -- (7.827,-1.577);
\draw[sym] (7.531,-1.651) -- (7.827,-1.651);
\node[lab,anchor=west] at (7.753,-1.391) {CTA};
\node[lab,anchor=west] at (7.753,-1.836) {\{HTH\}};
\draw[sym] (8.421,-1.985) -- (8.421,-2.059);
\draw[sym] (8.273,-1.910) -- (8.421,-1.985);
\draw[sym] (8.569,-1.762) -- (8.273,-1.910);
\draw[sym] (8.273,-1.614) -- (8.569,-1.762);
\draw[sym] (8.569,-1.465) -- (8.273,-1.614);
\draw[sym] (8.421,-1.317) -- (8.421,-1.391);
\draw[sym] (8.421,-1.391) -- (8.569,-1.465);
\node[lab,anchor=west] at (8.606,-1.540) {RTA};
\node[lab,anchor=west] at (8.606,-1.873) {\{1/GTH\}};
\draw[sym] (5.602,-3.469) -- (5.639,-3.580);
\draw[sym] (5.602,-3.469) -- (5.565,-3.580);
\draw[sym] (5.565,-3.580) -- (5.639,-3.580);
\draw[sym] (5.602,-3.580) -- (5.602,-3.765);
\draw[sym] (5.602,-3.988) -- (5.602,-3.914);
\draw[sym] (5.602,-3.246) -- (5.602,-3.320);
\draw[sym] (5.602,-3.617) ellipse (0.297 and 0.297);
\node[lab,anchor=west] at (5.824,-3.469) {BQB};
\node[lab,anchor=west] at (5.824,-3.765) {I=ETA*1e-6*V(PB)};
\draw[sym] (7.679,-3.580) -- (7.679,-3.840);
\draw[sym] (7.679,-3.506) -- (7.679,-3.246);
\draw[sym] (7.531,-3.506) -- (7.827,-3.506);
\draw[sym] (7.531,-3.580) -- (7.827,-3.580);
\node[lab,anchor=west] at (7.753,-3.320) {CTB};
\node[lab,anchor=west] at (7.753,-3.765) {\{HTH\}};
\draw[sym] (8.421,-3.914) -- (8.421,-3.988);
\draw[sym] (8.273,-3.840) -- (8.421,-3.914);
\draw[sym] (8.569,-3.691) -- (8.273,-3.840);
\draw[sym] (8.273,-3.543) -- (8.569,-3.691);
\draw[sym] (8.569,-3.394) -- (8.273,-3.543);
\draw[sym] (8.421,-3.246) -- (8.421,-3.320);
\draw[sym] (8.421,-3.320) -- (8.569,-3.394);
\node[lab,anchor=west] at (8.606,-3.469) {RTB};
\node[lab,anchor=west] at (8.606,-3.802) {\{1/GTH\}};
\draw[sym] (10.350,-1.836) -- (10.387,-1.948);
\draw[sym] (10.350,-1.836) -- (10.313,-1.948);
\draw[sym] (10.313,-1.948) -- (10.387,-1.948);
\draw[sym] (10.350,-1.948) -- (10.350,-2.133);
\draw[sym] (10.350,-2.356) -- (10.350,-2.281);
\draw[sym] (10.350,-1.614) -- (10.350,-1.688);
\draw[sym] (10.350,-1.985) ellipse (0.297 and 0.297);
\node[lab,anchor=west] at (10.573,-1.836) {BPYA};
\node[lab,anchor=west] at (10.573,-2.133) {I=PPY*AEL*ddt(V(TA))};
\draw[sym] (10.350,-2.875) -- (10.387,-2.764);
\draw[sym] (10.350,-2.875) -- (10.313,-2.764);
\draw[sym] (10.313,-2.764) -- (10.387,-2.764);
\draw[sym] (10.350,-2.578) -- (10.350,-2.764);
\draw[sym] (10.350,-3.098) -- (10.350,-3.023);
\draw[sym] (10.350,-2.356) -- (10.350,-2.430);
\draw[sym] (10.350,-2.727) ellipse (0.297 and 0.297);
\node[lab,anchor=west] at (10.573,-2.578) {BPYB};
\node[lab,anchor=west] at (10.573,-2.875) {I=PPY*(1-MISM)*AEL*ddt(V(TB))};
\draw[sym] (13.466,-1.948) -- (13.466,-2.207);
\draw[sym] (13.466,-1.873) -- (13.466,-1.614);
\draw[sym] (13.318,-1.873) -- (13.615,-1.873);
\draw[sym] (13.318,-1.948) -- (13.615,-1.948);
\node[lab,anchor=west] at (13.540,-1.688) {CEA};
\node[lab,anchor=west] at (13.540,-2.133) {\{CE\}};
\draw[sym] (13.466,-2.690) -- (13.466,-2.949);
\draw[sym] (13.466,-2.615) -- (13.466,-2.356);
\draw[sym] (13.318,-2.615) -- (13.615,-2.615);
\draw[sym] (13.318,-2.690) -- (13.615,-2.690);
\node[lab,anchor=west] at (13.540,-2.430) {CEB};
\node[lab,anchor=west] at (13.540,-2.875) {\{CE\}};
\draw[sym] (14.208,-2.281) -- (14.208,-2.356);
\draw[sym] (14.060,-2.207) -- (14.208,-2.281);
\draw[sym] (14.356,-2.059) -- (14.060,-2.207);
\draw[sym] (14.060,-1.910) -- (14.356,-2.059);
\draw[sym] (14.356,-1.762) -- (14.060,-1.910);
\draw[sym] (14.208,-1.614) -- (14.208,-1.688);
\draw[sym] (14.208,-1.688) -- (14.356,-1.762);
\node[lab,anchor=west] at (14.394,-1.836) {RLA};
\node[lab,anchor=west] at (14.394,-2.170) {1T};
\draw[sym] (14.208,-3.023) -- (14.208,-3.098);
\draw[sym] (14.060,-2.949) -- (14.208,-3.023);
\draw[sym] (14.356,-2.801) -- (14.060,-2.949);
\draw[sym] (14.060,-2.652) -- (14.356,-2.801);
\draw[sym] (14.356,-2.504) -- (14.060,-2.652);
\draw[sym] (14.208,-2.356) -- (14.208,-2.430);
\draw[sym] (14.208,-2.430) -- (14.356,-2.504);
\node[lab,anchor=west] at (14.394,-2.578) {RLB};
\node[lab,anchor=west] at (14.394,-2.912) {1T};
\draw[sym] (15.098,-2.281) -- (15.098,-2.356);
\draw[sym] (14.950,-2.207) -- (15.098,-2.281);
\draw[sym] (15.247,-2.059) -- (14.950,-2.207);
\draw[sym] (14.950,-1.910) -- (15.247,-2.059);
\draw[sym] (15.247,-1.762) -- (14.950,-1.910);
\draw[sym] (15.098,-1.614) -- (15.098,-1.688);
\draw[sym] (15.098,-1.688) -- (15.247,-1.762);
\node[lab,anchor=west] at (15.284,-1.836) {RG};
\node[lab,anchor=west] at (15.284,-2.170) {68G};
\draw[sym] (16.137,-1.169) -- (16.137,-1.762);
\draw[sym] (16.434,-1.688) -- (16.434,-1.910);
\draw[sym] (16.137,-1.688) -- (16.434,-1.688);
\draw[sym] (16.434,-1.243) -- (16.434,-1.020);
\draw[sym] (16.137,-1.243) -- (16.434,-1.243);
\draw[sym] (15.989,-1.614) -- (16.026,-1.614);
\draw[sym] (16.026,-1.651) -- (16.137,-1.614);
\draw[sym] (16.026,-1.577) -- (16.137,-1.614);
\draw[sym] (16.026,-1.577) -- (16.026,-1.651);
\node[lab,anchor=west] at (16.508,-1.317) {J1};
\node[lab,anchor=west] at (16.508,-1.688) {JLHI};
\draw[sym] (17.324,-2.244) -- (17.324,-2.504);
\draw[sym] (17.324,-2.170) -- (17.324,-1.910);
\draw[sym] (17.176,-2.170) -- (17.473,-2.170);
\draw[sym] (17.176,-2.244) -- (17.473,-2.244);
\node[lab,anchor=west] at (17.398,-1.985) {CEMI};
\node[lab,anchor=west] at (17.398,-2.430) {47p};
\draw[sym] (20.144,-2.578) -- (20.144,-2.652);
\draw[sym] (19.995,-2.504) -- (20.144,-2.578);
\draw[sym] (20.292,-2.356) -- (19.995,-2.504);
\draw[sym] (19.995,-2.207) -- (20.292,-2.356);
\draw[sym] (20.292,-2.059) -- (19.995,-2.207);
\draw[sym] (20.144,-1.910) -- (20.144,-1.985);
\draw[sym] (20.144,-1.985) -- (20.292,-2.059);
\node[lab,anchor=west] at (20.329,-2.133) {RS};
\node[lab,anchor=west] at (20.329,-2.467) {47k};
\draw[sym] (21.442,-1.836) -- (21.442,-1.985);
\draw[sym] (21.071,-1.836) -- (21.071,-1.985);
\draw[sym] (21.516,-1.910) -- (21.368,-1.910);
\draw[sym] (20.885,-1.910) -- (20.960,-1.910);
\draw[sym] (21.627,-1.910) -- (21.553,-1.910);
\draw[sym] (21.256,-1.910) ellipse (0.297 and 0.297);
\node[lab,anchor=south] at (21.256,-1.614) {BN};
\node[lab,anchor=north] at (21.256,-2.207) {V=VNOISE(time)};
\end{tikzpicture}

\caption{Tín hiệu vào, mô hình nhiệt, mô hình \LHI{} (khung nét đứt), $R_S$ và nguồn nhiễu BN}
\end{subfigure}

\vspace{0.3cm}
\begin{subfigure}{\linewidth}\centering
\begin{tikzpicture}[x=1cm,y=1cm,line cap=round,line join=round,wire/.style={line width=0.5pt,c1!75!black},sym/.style={line width=0.45pt,bkred!75!black},lab/.style={font=\fontsize{4.6}{5}\selectfont\sffamily,inner sep=0.5pt},net/.style={font=\fontsize{4.8}{5}\selectfont\sffamily\bfseries,text=c2!85!black,inner sep=0.6pt}]
\draw[wire] (6.080,-2.934) -- (6.080,-2.511);
\draw[wire] (6.080,-3.780) -- (6.080,-4.203);
\draw[wire] (5.657,-3.569) -- (4.600,-3.569);
\draw[wire] (5.657,-3.146) -- (5.234,-3.146);
\draw[wire] (5.234,-3.146) -- (5.234,-1.877);
\draw[wire] (5.023,-1.877) -- (5.234,-1.877);
\draw[wire] (5.234,-1.877) -- (5.657,-1.877);
\draw[wire] (5.234,-1.877) -- (5.234,-0.820);
\draw[wire] (5.234,-0.820) -- (5.869,-0.820);
\draw[wire] (6.715,-1.877) -- (7.349,-1.877);
\draw[wire] (6.715,-0.820) -- (7.349,-0.820);
\draw[wire] (7.349,-0.820) -- (7.349,-1.877);
\draw[wire] (7.349,-1.877) -- (7.349,-3.357);
\draw[wire] (6.503,-3.357) -- (7.349,-3.357);
\draw[wire] (3.966,-1.877) -- (3.331,-1.877);
\draw[wire] (11.156,-2.723) -- (11.156,-2.300);
\draw[wire] (11.156,-3.569) -- (11.156,-3.992);
\draw[wire] (7.349,-3.357) -- (8.195,-3.357);
\draw[wire] (9.041,-3.357) -- (9.887,-3.357);
\draw[wire] (9.887,-3.357) -- (10.733,-3.357);
\draw[wire] (9.887,-4.415) -- (9.887,-4.838);
\draw[wire] (10.733,-2.934) -- (10.310,-2.934);
\draw[wire] (10.310,-2.934) -- (10.310,-1.666);
\draw[wire] (10.099,-1.666) -- (10.310,-1.666);
\draw[wire] (10.310,-1.666) -- (10.733,-1.666);
\draw[wire] (10.310,-1.666) -- (10.310,-0.608);
\draw[wire] (10.310,-0.608) -- (10.945,-0.608);
\draw[wire] (11.791,-1.666) -- (12.425,-1.666);
\draw[wire] (11.791,-0.608) -- (12.425,-0.608);
\draw[wire] (12.425,-0.608) -- (12.425,-1.666);
\draw[wire] (12.425,-1.666) -- (12.425,-3.146);
\draw[wire] (11.579,-3.146) -- (12.425,-3.146);
\draw[wire] (9.041,-1.666) -- (8.618,-1.666);
\draw[wire] (17.713,-2.511) -- (17.713,-2.089);
\draw[wire] (17.713,-3.357) -- (17.713,-3.780);
\draw[wire] (12.425,-3.146) -- (12.848,-3.146);
\draw[wire] (13.906,-3.146) -- (14.329,-3.146);
\draw[wire] (14.329,-3.146) -- (14.540,-3.146);
\draw[wire] (15.598,-3.146) -- (16.021,-3.146);
\draw[wire] (16.021,-3.146) -- (17.290,-3.146);
\draw[wire] (14.329,-1.666) -- (14.329,-3.146);
\draw[wire] (16.021,-3.992) -- (16.021,-4.415);
\draw[wire] (14.329,-0.820) -- (18.982,-0.820);
\draw[wire] (18.982,-0.820) -- (18.982,-1.243);
\draw[wire] (18.982,-1.243) -- (18.982,-2.934);
\draw[wire] (17.290,-2.723) -- (16.867,-2.723);
\draw[wire] (16.867,-2.723) -- (16.867,-1.243);
\draw[wire] (16.867,-1.243) -- (18.982,-1.243);
\draw[wire] (18.136,-2.934) -- (18.982,-2.934);
\draw[wire] (18.982,-2.934) -- (19.405,-2.934);
\draw[wire] (20.462,-2.934) -- (21.308,-2.934);
\draw[wire] (21.308,-2.934) -- (22.366,-2.934);
\fill[c2!85!black] (6.080,-2.511) circle (0.6pt);
\node[net,anchor=south west] at (6.080,-2.511) {VDD};
\draw[wire] (6.080,-4.203) -- ++(0,-0.106) ++(-0.211,0) -- ++(0.423,0) -- ++(-0.211,-0.211) -- cycle;
\fill[c2!85!black] (4.600,-3.569) circle (0.6pt);
\node[net,anchor=south west] at (4.600,-3.569) {SN};
\draw[wire] (3.331,-2.723) -- ++(0,-0.106) ++(-0.211,0) -- ++(0.423,0) -- ++(-0.211,-0.211) -- cycle;
\fill[c2!85!black] (7.349,-3.357) circle (0.6pt);
\node[net,anchor=south west] at (7.349,-3.357) {A1};
\fill[c2!85!black] (11.156,-2.300) circle (0.6pt);
\node[net,anchor=south west] at (11.156,-2.300) {VDD};
\draw[wire] (11.156,-3.992) -- ++(0,-0.106) ++(-0.211,0) -- ++(0.423,0) -- ++(-0.211,-0.211) -- cycle;
\fill[c2!85!black] (9.887,-4.838) circle (0.6pt);
\node[net,anchor=south west] at (9.887,-4.838) {VREF};
\fill[c2!85!black] (8.618,-1.666) circle (0.6pt);
\node[net,anchor=south west] at (8.618,-1.666) {VREF};
\fill[c2!85!black] (12.425,-3.146) circle (0.6pt);
\node[net,anchor=south west] at (12.425,-3.146) {A2};
\fill[c2!85!black] (17.713,-2.089) circle (0.6pt);
\node[net,anchor=south west] at (17.713,-2.089) {VDD};
\draw[wire] (17.713,-3.780) -- ++(0,-0.106) ++(-0.211,0) -- ++(0.423,0) -- ++(-0.211,-0.211) -- cycle;
\fill[c2!85!black] (16.021,-4.415) circle (0.6pt);
\node[net,anchor=south west] at (16.021,-4.415) {VREF};
\fill[c2!85!black] (18.982,-2.934) circle (0.6pt);
\node[net,anchor=south west] at (18.982,-2.934) {A3};
\draw[wire] (21.308,-3.992) -- ++(0,-0.106) ++(-0.211,0) -- ++(0.423,0) -- ++(-0.211,-0.211) -- cycle;
\fill[c2!85!black] (22.366,-2.934) circle (0.6pt);
\node[net,anchor=south west] at (22.366,-2.934) {OUT};
\draw[sym] (5.657,-2.934) -- (6.503,-3.357);
\draw[sym] (5.657,-3.780) -- (6.503,-3.357);
\draw[sym] (5.657,-2.934) -- (5.657,-3.780);
\draw[sym] (5.710,-3.146) -- (5.816,-3.146);
\draw[sym] (5.710,-3.569) -- (5.816,-3.569);
\draw[sym] (5.763,-3.622) -- (5.763,-3.516);
\draw[sym] (6.080,-2.934) -- (6.080,-3.146);
\draw[sym] (6.080,-3.780) -- (6.080,-3.569);
\draw[sym] (6.133,-3.093) -- (6.239,-3.093);
\draw[sym] (6.186,-3.040) -- (6.186,-3.146);
\draw[sym] (6.133,-3.622) -- (6.239,-3.622);
\node[lab,anchor=west] at (6.292,-2.934) {U1A};
\node[lab,anchor=west] at (6.292,-3.780) {AD8608};
\draw[sym] (5.763,-1.877) -- (5.657,-1.877);
\draw[sym] (5.869,-1.666) -- (5.763,-1.877);
\draw[sym] (6.080,-2.089) -- (5.869,-1.666);
\draw[sym] (6.292,-1.666) -- (6.080,-2.089);
\draw[sym] (6.503,-2.089) -- (6.292,-1.666);
\draw[sym] (6.715,-1.877) -- (6.609,-1.877);
\draw[sym] (6.609,-1.877) -- (6.503,-2.089);
\node[lab,anchor=south] at (6.186,-1.666) {R2};
\node[lab,anchor=north] at (6.186,-2.089) {1Meg};
\draw[sym] (6.239,-0.820) -- (5.869,-0.820);
\draw[sym] (6.345,-0.820) -- (6.715,-0.820);
\draw[sym] (6.345,-0.608) -- (6.345,-1.031);
\draw[sym] (6.239,-0.608) -- (6.239,-1.031);
\node[lab,anchor=south] at (6.292,-0.608) {C2};
\node[lab,anchor=north] at (6.292,-1.031) {10n};
\draw[sym] (4.071,-1.877) -- (3.966,-1.877);
\draw[sym] (4.177,-1.666) -- (4.071,-1.877);
\draw[sym] (4.389,-2.089) -- (4.177,-1.666);
\draw[sym] (4.600,-1.666) -- (4.389,-2.089);
\draw[sym] (4.811,-2.089) -- (4.600,-1.666);
\draw[sym] (5.023,-1.877) -- (4.917,-1.877);
\draw[sym] (4.917,-1.877) -- (4.811,-2.089);
\node[lab,anchor=south] at (4.494,-1.666) {R1};
\node[lab,anchor=north] at (4.494,-2.089) {10k};
\draw[sym] (3.331,-2.353) -- (3.331,-2.723);
\draw[sym] (3.331,-1.877) -- (3.331,-2.247);
\draw[sym] (3.225,-2.036) -- (3.225,-2.141);
\draw[sym] (3.172,-2.089) -- (3.278,-2.089);
\draw[sym] (3.120,-2.247) -- (3.543,-2.247);
\draw[sym] (3.541,-2.409) arc[start angle=60.3,end angle=119.7,x radius=0.423,y radius=0.423];
\node[lab,anchor=west] at (3.437,-1.983) {C1};
\node[lab,anchor=west] at (3.437,-2.630) {47u};
\draw[sym] (8.566,-3.357) -- (8.195,-3.357);
\draw[sym] (8.671,-3.357) -- (9.041,-3.357);
\draw[sym] (8.671,-3.146) -- (8.671,-3.569);
\draw[sym] (8.566,-3.146) -- (8.566,-3.569);
\node[lab,anchor=south] at (8.618,-3.146) {C3};
\node[lab,anchor=north] at (8.618,-3.569) {1u};
\draw[sym] (9.887,-4.309) -- (9.887,-4.415);
\draw[sym] (9.676,-4.203) -- (9.887,-4.309);
\draw[sym] (10.099,-3.992) -- (9.676,-4.203);
\draw[sym] (9.676,-3.780) -- (10.099,-3.992);
\draw[sym] (10.099,-3.569) -- (9.676,-3.780);
\draw[sym] (9.887,-3.357) -- (9.887,-3.463);
\draw[sym] (9.887,-3.463) -- (10.099,-3.569);
\node[lab,anchor=west] at (10.152,-3.675) {R3};
\node[lab,anchor=west] at (10.152,-4.151) {470k};
\draw[sym] (10.733,-2.723) -- (11.579,-3.146);
\draw[sym] (10.733,-3.569) -- (11.579,-3.146);
\draw[sym] (10.733,-2.723) -- (10.733,-3.569);
\draw[sym] (10.786,-2.934) -- (10.892,-2.934);
\draw[sym] (10.786,-3.357) -- (10.892,-3.357);
\draw[sym] (10.839,-3.410) -- (10.839,-3.305);
\draw[sym] (11.156,-2.723) -- (11.156,-2.934);
\draw[sym] (11.156,-3.569) -- (11.156,-3.357);
\draw[sym] (11.209,-2.882) -- (11.315,-2.882);
\draw[sym] (11.262,-2.829) -- (11.262,-2.934);
\draw[sym] (11.209,-3.410) -- (11.315,-3.410);
\node[lab,anchor=west] at (11.368,-2.723) {U1B};
\node[lab,anchor=west] at (11.368,-3.569) {AD8608};
\draw[sym] (10.839,-1.666) -- (10.733,-1.666);
\draw[sym] (10.945,-1.454) -- (10.839,-1.666);
\draw[sym] (11.156,-1.877) -- (10.945,-1.454);
\draw[sym] (11.368,-1.454) -- (11.156,-1.877);
\draw[sym] (11.579,-1.877) -- (11.368,-1.454);
\draw[sym] (11.791,-1.666) -- (11.685,-1.666);
\draw[sym] (11.685,-1.666) -- (11.579,-1.877);
\node[lab,anchor=south] at (11.262,-1.454) {R5};
\node[lab,anchor=north] at (11.262,-1.877) {150k};
\draw[sym] (11.315,-0.608) -- (10.945,-0.608);
\draw[sym] (11.421,-0.608) -- (11.791,-0.608);
\draw[sym] (11.421,-0.397) -- (11.421,-0.820);
\draw[sym] (11.315,-0.397) -- (11.315,-0.820);
\node[lab,anchor=south] at (11.368,-0.397) {C4};
\node[lab,anchor=north] at (11.368,-0.820) {68n};
\draw[sym] (9.147,-1.666) -- (9.041,-1.666);
\draw[sym] (9.253,-1.454) -- (9.147,-1.666);
\draw[sym] (9.464,-1.877) -- (9.253,-1.454);
\draw[sym] (9.676,-1.454) -- (9.464,-1.877);
\draw[sym] (9.887,-1.877) -- (9.676,-1.454);
\draw[sym] (10.099,-1.666) -- (9.993,-1.666);
\draw[sym] (9.993,-1.666) -- (9.887,-1.877);
\node[lab,anchor=south] at (9.570,-1.454) {R4};
\node[lab,anchor=north] at (9.570,-1.877) {10k};
\draw[sym] (17.290,-2.511) -- (18.136,-2.934);
\draw[sym] (17.290,-3.357) -- (18.136,-2.934);
\draw[sym] (17.290,-2.511) -- (17.290,-3.357);
\draw[sym] (17.343,-2.723) -- (17.448,-2.723);
\draw[sym] (17.343,-3.146) -- (17.448,-3.146);
\draw[sym] (17.395,-3.199) -- (17.395,-3.093);
\draw[sym] (17.713,-2.511) -- (17.713,-2.723);
\draw[sym] (17.713,-3.357) -- (17.713,-3.146);
\draw[sym] (17.766,-2.670) -- (17.871,-2.670);
\draw[sym] (17.818,-2.617) -- (17.818,-2.723);
\draw[sym] (17.766,-3.199) -- (17.871,-3.199);
\node[lab,anchor=west] at (17.924,-2.511) {U1C};
\node[lab,anchor=west] at (17.924,-3.357) {AD8608};
\draw[sym] (12.954,-3.146) -- (12.848,-3.146);
\draw[sym] (13.060,-2.934) -- (12.954,-3.146);
\draw[sym] (13.271,-3.357) -- (13.060,-2.934);
\draw[sym] (13.483,-2.934) -- (13.271,-3.357);
\draw[sym] (13.694,-3.357) -- (13.483,-2.934);
\draw[sym] (13.906,-3.146) -- (13.800,-3.146);
\draw[sym] (13.800,-3.146) -- (13.694,-3.357);
\node[lab,anchor=south] at (13.377,-2.934) {R6};
\node[lab,anchor=north] at (13.377,-3.357) {100k};
\draw[sym] (14.646,-3.146) -- (14.540,-3.146);
\draw[sym] (14.752,-2.934) -- (14.646,-3.146);
\draw[sym] (14.963,-3.357) -- (14.752,-2.934);
\draw[sym] (15.175,-2.934) -- (14.963,-3.357);
\draw[sym] (15.386,-3.357) -- (15.175,-2.934);
\draw[sym] (15.598,-3.146) -- (15.492,-3.146);
\draw[sym] (15.492,-3.146) -- (15.386,-3.357);
\node[lab,anchor=south] at (15.069,-2.934) {R7};
\node[lab,anchor=north] at (15.069,-3.357) {100k};
\draw[sym] (14.329,-1.295) -- (14.329,-1.666);
\draw[sym] (14.329,-1.190) -- (14.329,-0.820);
\draw[sym] (14.117,-1.190) -- (14.540,-1.190);
\draw[sym] (14.117,-1.295) -- (14.540,-1.295);
\node[lab,anchor=west] at (14.434,-0.925) {C5};
\node[lab,anchor=west] at (14.434,-1.560) {220n};
\draw[sym] (16.021,-3.622) -- (16.021,-3.992);
\draw[sym] (16.021,-3.516) -- (16.021,-3.146);
\draw[sym] (15.809,-3.516) -- (16.232,-3.516);
\draw[sym] (15.809,-3.622) -- (16.232,-3.622);
\node[lab,anchor=west] at (16.126,-3.252) {C6};
\node[lab,anchor=west] at (16.126,-3.886) {100n};
\draw[sym] (19.510,-2.934) -- (19.405,-2.934);
\draw[sym] (19.616,-2.723) -- (19.510,-2.934);
\draw[sym] (19.828,-3.146) -- (19.616,-2.723);
\draw[sym] (20.039,-2.723) -- (19.828,-3.146);
\draw[sym] (20.251,-3.146) -- (20.039,-2.723);
\draw[sym] (20.462,-2.934) -- (20.356,-2.934);
\draw[sym] (20.356,-2.934) -- (20.251,-3.146);
\node[lab,anchor=south] at (19.933,-2.723) {R8};
\node[lab,anchor=north] at (19.933,-3.146) {10k};
\draw[sym] (21.308,-3.886) -- (21.308,-3.992);
\draw[sym] (21.097,-3.780) -- (21.308,-3.886);
\draw[sym] (21.520,-3.569) -- (21.097,-3.780);
\draw[sym] (21.097,-3.357) -- (21.520,-3.569);
\draw[sym] (21.520,-3.146) -- (21.097,-3.357);
\draw[sym] (21.308,-2.934) -- (21.308,-3.040);
\draw[sym] (21.308,-3.040) -- (21.520,-3.146);
\node[lab,anchor=west] at (21.572,-3.252) {R9};
\node[lab,anchor=west] at (21.572,-3.728) {100k};
\end{tikzpicture}

\caption{Mạch khuếch đại (U1A, U1B), lọc Sallen-Key (U1C) và chia áp đầu ra $R_8$ -- $R_9$}
\end{subfigure}

\vspace{0.3cm}
\begin{subfigure}{\linewidth}\centering
\begin{tikzpicture}[x=1cm,y=1cm,line cap=round,line join=round,wire/.style={line width=0.5pt,c1!75!black},sym/.style={line width=0.45pt,bkred!75!black},lab/.style={font=\fontsize{4.6}{5}\selectfont\sffamily,inner sep=0.5pt},net/.style={font=\fontsize{4.8}{5}\selectfont\sffamily\bfseries,text=c2!85!black,inner sep=0.6pt}]
\draw[wire] (4.811,-2.591) -- (6.080,-2.591);
\draw[wire] (6.080,-2.591) -- (7.349,-2.591);
\draw[wire] (7.772,-1.956) -- (7.772,-1.533);
\draw[wire] (7.772,-2.802) -- (7.772,-3.225);
\draw[wire] (7.349,-2.168) -- (6.926,-2.168);
\draw[wire] (6.926,-2.168) -- (6.926,-0.687);
\draw[wire] (6.926,-0.687) -- (9.041,-0.687);
\draw[wire] (9.041,-0.687) -- (9.041,-2.379);
\draw[wire] (8.195,-2.379) -- (9.041,-2.379);
\draw[wire] (9.041,-2.379) -- (9.887,-2.379);
\fill[c2!85!black] (2.062,-1.533) circle (0.6pt);
\node[net,anchor=south west] at (2.062,-1.533) {VDD};
\draw[wire] (2.062,-2.591) -- ++(0,-0.106) ++(-0.211,0) -- ++(0.423,0) -- ++(-0.211,-0.211) -- cycle;
\fill[c2!85!black] (4.811,-1.533) circle (0.6pt);
\node[net,anchor=south west] at (4.811,-1.533) {VDD};
\draw[wire] (4.811,-3.648) -- ++(0,-0.106) ++(-0.211,0) -- ++(0.423,0) -- ++(-0.211,-0.211) -- cycle;
\draw[wire] (6.080,-3.437) -- ++(0,-0.106) ++(-0.211,0) -- ++(0.423,0) -- ++(-0.211,-0.211) -- cycle;
\fill[c2!85!black] (7.772,-1.533) circle (0.6pt);
\node[net,anchor=south west] at (7.772,-1.533) {VDD};
\draw[wire] (7.772,-3.225) -- ++(0,-0.106) ++(-0.211,0) -- ++(0.423,0) -- ++(-0.211,-0.211) -- cycle;
\fill[c2!85!black] (9.887,-2.379) circle (0.6pt);
\node[net,anchor=south west] at (9.887,-2.379) {VREF};
\draw[sym] (1.956,-1.798) -- (2.168,-1.798);
\draw[sym] (1.956,-2.326) -- (2.168,-2.326);
\draw[sym] (2.062,-1.692) -- (2.062,-1.903);
\draw[sym] (2.062,-2.591) -- (2.062,-2.485);
\draw[sym] (2.062,-1.533) -- (2.062,-1.639);
\draw[sym] (2.062,-2.062) ellipse (0.423 and 0.423);
\node[lab,anchor=west] at (2.379,-1.533) {V1};
\node[lab,anchor=west] at (2.379,-2.591) {3.3};
\draw[sym] (4.811,-2.485) -- (4.811,-2.591);
\draw[sym] (4.600,-2.379) -- (4.811,-2.485);
\draw[sym] (5.023,-2.168) -- (4.600,-2.379);
\draw[sym] (4.600,-1.956) -- (5.023,-2.168);
\draw[sym] (5.023,-1.745) -- (4.600,-1.956);
\draw[sym] (4.811,-1.533) -- (4.811,-1.639);
\draw[sym] (4.811,-1.639) -- (5.023,-1.745);
\node[lab,anchor=west] at (5.076,-1.851) {R10};
\node[lab,anchor=west] at (5.076,-2.326) {100k};
\draw[sym] (4.811,-3.542) -- (4.811,-3.648);
\draw[sym] (4.600,-3.437) -- (4.811,-3.542);
\draw[sym] (5.023,-3.225) -- (4.600,-3.437);
\draw[sym] (4.600,-3.014) -- (5.023,-3.225);
\draw[sym] (5.023,-2.802) -- (4.600,-3.014);
\draw[sym] (4.811,-2.591) -- (4.811,-2.696);
\draw[sym] (4.811,-2.696) -- (5.023,-2.802);
\node[lab,anchor=west] at (5.076,-2.908) {R11};
\node[lab,anchor=west] at (5.076,-3.384) {100k};
\draw[sym] (6.080,-3.067) -- (6.080,-3.437);
\draw[sym] (6.080,-2.961) -- (6.080,-2.591);
\draw[sym] (5.869,-2.961) -- (6.292,-2.961);
\draw[sym] (5.869,-3.067) -- (6.292,-3.067);
\node[lab,anchor=west] at (6.186,-2.696) {C7};
\node[lab,anchor=west] at (6.186,-3.331) {10u};
\draw[sym] (7.349,-1.956) -- (8.195,-2.379);
\draw[sym] (7.349,-2.802) -- (8.195,-2.379);
\draw[sym] (7.349,-1.956) -- (7.349,-2.802);
\draw[sym] (7.402,-2.168) -- (7.508,-2.168);
\draw[sym] (7.402,-2.591) -- (7.508,-2.591);
\draw[sym] (7.455,-2.644) -- (7.455,-2.538);
\draw[sym] (7.772,-1.956) -- (7.772,-2.168);
\draw[sym] (7.772,-2.802) -- (7.772,-2.591);
\draw[sym] (7.825,-2.115) -- (7.931,-2.115);
\draw[sym] (7.878,-2.062) -- (7.878,-2.168);
\draw[sym] (7.825,-2.644) -- (7.931,-2.644);
\node[lab,anchor=west] at (7.984,-1.956) {U1D};
\node[lab,anchor=west] at (7.984,-2.802) {AD8608};
\end{tikzpicture}

\caption{Nguồn 3,3~V và mạch tạo $\VREF$ (U1D)}
\end{subfigure}
\caption{Sơ đồ nguyên lý toàn hệ thống trong file LTspice \texttt{PIR\_LHI778.asc} (vẽ lại từ file mô phỏng)}
\label{fig:ltspice}
\end{figure}
\end{landscape}

\subsection{Điểm làm việc tĩnh}

Khi không có chuyển động (0,5 -- 2,5~s), các nút có điện áp như trong bảng~\ref{tab:dc}.

\begin{table}[!htb]
\centering
\caption{Điểm làm việc tĩnh (không có chuyển động)}
\label{tab:dc}
\small
\begin{tabularx}{\textwidth}{@{}p{3.3cm}p{2.6cm}p{2.8cm}L@{}}
\toprule
\textbf{Nút} & \textbf{Mô phỏng} & \textbf{Tính toán} & \textbf{Nhận xét} \\ \midrule
Cổng FET $V_G$ & 0,4~mV & $\approx 0$ & Cổng nối đất qua $R_G$ \\
Chân S, $V_S$ & 0,8095~V & 0,805~V & Trong dải datasheet 0,2 -- 1,55~V \\
Tầng 1, $V_{A1}$ & 0,8094~V & $= V_S$ & $C_1$ chặn DC, hệ số DC bằng 1 \\
$\VREF$ & 1,6500~V & 1,65~V & $V_{DD}/2$ \\
Tầng 2, $V_{A2}$ & 1,6492~V & $\VREF \pm 1{,}2$~mV & Lệch $-0{,}8$~mV do $16\,V_{OS}$ \\
Bộ lọc, $V_{A3}$ & 1,6491~V & $= V_{A2}$ & Hệ số 1 \\
\textbf{Đầu ra OUT} & \textbf{1,4992~V} & \textbf{1,50~V} & \textbf{Đạt yêu cầu 1,5~V (sai số 0,05\%)} \\
\bottomrule
\end{tabularx}
\end{table}

Mức nghỉ của đầu ra là 1,4992~V, chỉ lệch 0,8~mV so với 1,5~V yêu cầu. Nhờ tụ $C_1$ và $C_3$, điện áp offset 0,81~V của cảm biến \emph{không} ảnh hưởng tới mức nghỉ. Nếu thay bằng một cảm biến khác có $V_S = 0{,}2$~V hay 1,55~V, đầu ra vẫn ở 1,5~V.

\subsection{Đáp ứng tần số của mạch chuẩn hóa}

Hình~\ref{fig:bode} trình bày hệ số khuếch đại của từng khối và của toàn mạch (từ chân S tới OUT), lấy từ phân tích AC. Các số liệu chính tổng hợp trong bảng~\ref{tab:bode}.

\begin{figure}[!htb]
\centering
\begin{tikzpicture}
\begin{semilogxaxis}[width=0.93\textwidth,height=8.2cm,xmin=0.01,xmax=1000,ymin=-90,ymax=70,
  xlabel={Tần số [Hz]},ylabel={Hệ số khuếch đại [dB]},ytick distance=20,
  legend pos=south west,log ticks with fixed point]
  \fill[c3!10] (axis cs:0.5,-90) rectangle (axis cs:8.5,70);
  \addplot[c2,thick] table[x=f,y=G1] {data/ac.dat}; \addlegendentry{Tầng 1: $V_{A1}/V_S$}
  \addplot[c4,thick] table[x=f,y=G2] {data/ac.dat}; \addlegendentry{Ghép AC + tầng 2: $V_{A2}/V_{A1}$}
  \addplot[c3,thick,dashed] table[x=f,y=G3] {data/ac.dat}; \addlegendentry{Sallen-Key: $V_{A3}/V_{A2}$}
  \addplot[c1,very thick] table[x=f,y=G] {data/ac.dat}; \addlegendentry{Toàn mạch: $V_\text{OUT}/V_S$}
  \draw[densely dashed,cgray] (axis cs:0.01,60) -- (axis cs:1000,60) node[pos=0.97,above,font=\scriptsize]{$63-3=60$ dB};
  \draw[densely dotted] (axis cs:50,-90) -- (axis cs:50,70) node[pos=0.95,left,font=\scriptsize]{50 Hz};
  \node[font=\scriptsize,c3!60!black] at (axis cs:2.1,-20) {dải thông 0,5 -- 8,5 Hz};
\end{semilogxaxis}
\end{tikzpicture}
\caption{Đáp ứng tần số của các tầng và của toàn mạch chuẩn hóa (mô phỏng AC)}
\label{fig:bode}
\end{figure}

\begin{table}[!htb]
\centering
\caption{Hệ số khuếch đại của từng khối theo tần số (mô phỏng AC)}
\label{tab:bode}
\small
\begin{tabularx}{\textwidth}{@{}p{3.2cm}YYYYY>{\centering\arraybackslash}p{3.0cm}@{}}
\toprule
\textbf{Khối} & \textbf{1 Hz} & \textbf{10 Hz} & \textbf{50 Hz} & \textbf{100 Hz} & \textbf{1 kHz} & \textbf{Dải $-3$ dB} \\ \midrule
Tầng 1 & 39,6 & 38,6 & 29,7 & 24,0 & 5,5 & 0,33 -- 16,2 Hz \\
Ghép AC + tầng 2 & 23,6 & 22,6 & 13,7 & 8,5 & 0,3 & 0,33 -- 16,2 Hz \\
Sallen-Key & 0,0 & $-2{,}0$ & $-26{,}8$ & $-38{,}8$ & $-78{,}7$ & $<11$ Hz \\
\textbf{Toàn mạch} [dB] & \textbf{62,4} & \textbf{58,4} & \textbf{15,9} & $\mathbf{-7{,}1}$ & $\mathbf{-73{,}8}$ & \textbf{0,50 -- 8,5 Hz} \\
\bottomrule
\end{tabularx}
\end{table}

Các kết quả mô phỏng khớp với tính toán ở mục~\ref{sec:thietke}:
\begin{itemize}
  \item Hệ số khuếch đại cực đại 63,0~dB ($\approx 1413$ lần), sát giá trị thiết kế 63,3~dB. Tại 1~Hz đạt 62,4~dB (1314 lần).
  \item Tần số cắt dưới 0,50~Hz (tính toán 0,53~Hz) và tần số cắt trên 8,5~Hz, đúng như dải thông mong muốn.
  \item So với dải thông, nhiễu 50~Hz bị suy giảm $63{,}0 - 15{,}9 = 47$~dB (khoảng 225 lần). Nhiễu 1~kHz bị suy giảm 137~dB.
\end{itemize}

\subsection{Tín hiệu bên trong cảm biến}
\label{sec:kq-cambien}

Hình~\ref{fig:sk1chain} theo dõi tín hiệu qua từng bước biến đổi bên trong \LHI{} ở kịch bản SK1 (người đi bình thường). Hình này minh họa trực tiếp nguyên lý đo ở mục~\ref{sec:hieuung} -- \ref{sec:cautao}.

\begin{figure}[!htb]
\centering
\begin{tikzpicture}
\begin{groupplot}[group style={group size=1 by 4,vertical sep=0.25cm,xticklabels at=edge bottom,xlabels at=edge bottom},
  width=0.93\textwidth,height=3.3cm,xmin=2.5,xmax=7.5,xlabel={Thời gian [s]},
  legend style={at={(1.0,1.0)},anchor=north east,legend columns=2}]
  \nextgroupplot[ylabel={$P$ [µW]},ymin=-0.02,ymax=0.16]
    \addplot[c1,thick] table[x=t,y=pa] {data/tran_sk1.dat}; \addlegendentry{$P_A$}
    \addplot[c2,thick] table[x=t,y=pb] {data/tran_sk1.dat}; \addlegendentry{$P_B$}
  \nextgroupplot[ylabel={$\Delta T$ [mK]},ymin=-0.01,ymax=0.12]
    \addplot[c1,thick] table[x=t,y=ta] {data/tran_sk1.dat}; \addlegendentry{$T_A$}
    \addplot[c2,thick] table[x=t,y=tb] {data/tran_sk1.dat}; \addlegendentry{$T_B$}
  \nextgroupplot[ylabel={$i_p$ [pA]},ymin=-0.3,ymax=0.3]
    \addplot[c1,thick] table[x=t,y=ia] {data/tran_sk1.dat}; \addlegendentry{$i_A$}
    \addplot[c2,thick] table[x=t,y=ib] {data/tran_sk1.dat}; \addlegendentry{$i_B$}
  \nextgroupplot[ylabel={$\Delta V_S$ [mV]},ymin=-0.7,ymax=0.7]
    \addplot[c3,very thick] table[x=t,y=s] {data/tran_sk1.dat}; \addlegendentry{chân S}
\end{groupplot}
\end{tikzpicture}
\caption{Kịch bản SK1: công suất tới, nhiệt độ phần tử, dòng nhiệt điện và điện áp chân S ($\Delta V_S = V_S - V_{S0}$, $V_{S0} = 0{,}8095$~V)}
\label{fig:sk1chain}
\end{figure}

Nhận xét:
\begin{itemize}
  \item \textbf{Nhiệt độ}: với công suất 0,13~µW, phần tử chỉ nóng lên tối đa \textbf{0,098~mK}. Đường $T(t)$ trễ và tròn hơn xung công suất do quán tính nhiệt $\tau_T = 0{,}6$~s.
  \item \textbf{Dòng nhiệt điện}: $i_p$ tỉ lệ với \emph{đạo hàm} nhiệt độ. Nó dương khi phần tử đang nóng lên, âm khi đang nguội đi, và có đỉnh chỉ khoảng \textbf{0,23~pA}. Đây là minh chứng trực tiếp cho công thức \eqref{eq:ip}.
  \item \textbf{Điện áp chân S}: dòng của hai phần tử đi vào cổng FET với dấu ngược nhau, nên $V_S$ là một dao động dương -- âm xen kẽ với biên độ \textbf{1,03~mV$_\text{pp}$} quanh $V_{S0} = 0{,}8095$~V. Biên độ này phù hợp với mức ``khoảng 1~mV$_\text{pp}$'' của tín hiệu PIR thực tế \cite{st-an4368}.
\end{itemize}

\subsection{Tín hiệu qua các tầng của mạch chuẩn hóa}
\label{sec:kq-tang}

Hình~\ref{fig:sk2stages} theo dõi kịch bản SK2 (tín hiệu lớn nhất) qua từng tầng của mạch. Bảng~\ref{tab:bientang} tổng hợp biên độ đỉnh -- đỉnh tại mỗi tầng cho ba kịch bản người đi.

\begin{figure}[!htb]
\centering
\begin{tikzpicture}
\begin{groupplot}[group style={group size=1 by 4,vertical sep=0.25cm,xticklabels at=edge bottom,xlabels at=edge bottom},
  width=0.93\textwidth,height=3.3cm,xmin=9.5,xmax=13.5,xlabel={Thời gian [s]},
  legend style={at={(1.0,1.0)},anchor=north east,legend columns=2}]
  \nextgroupplot[ylabel={$\Delta V_S$ [mV]},ymin=-1.3,ymax=1.3]
    \addplot[c3,thick] table[x=t,y=s] {data/tran_sk2.dat}; \addlegendentry{chân S ($\approx$ mV)}
  \nextgroupplot[ylabel={$V_{A1}$ [V]},ymin=0.68,ymax=0.94]
    \addplot[c2,thick] table[x=t,y=a1] {data/tran_sk2.dat}; \addlegendentry{ra tầng 1}
  \nextgroupplot[ylabel={[V]},ymin=0,ymax=3.3]
    \addplot[c4,thin] table[x=t,y=a2] {data/tran_sk2.dat}; \addlegendentry{$V_{A2}$ (trước lọc)}
    \addplot[c1,thick] table[x=t,y=a3] {data/tran_sk2.dat}; \addlegendentry{$V_{A3}$ (sau lọc)}
  \nextgroupplot[ylabel={$V_\text{OUT}$ [V]},ymin=0,ymax=3.3,ytick={0,0.5,...,3}]
    \addplot[c1,very thick] table[x=t,y=out] {data/tran_sk2.dat}; \addlegendentry{OUT}
    \draw[bkred,densely dashed] (axis cs:9.5,3) -- (axis cs:13.5,3);
    \draw[cgray,densely dashed] (axis cs:9.5,1.5) -- (axis cs:13.5,1.5);
\end{groupplot}
\end{tikzpicture}
\caption{Kịch bản SK2: tín hiệu tại chân S và sau từng tầng của mạch}
\label{fig:sk2stages}
\end{figure}

\begin{table}[!htb]
\centering
\caption{Biên độ đỉnh -- đỉnh tại các tầng (mô phỏng quá độ)}
\label{tab:bientang}
\small
\begin{tabularx}{\textwidth}{@{}p{2.5cm}YYYYY@{}}
\toprule
\textbf{Kịch bản} & $V_S$ [mV$_\text{pp}$] & $V_{A1}$ [V$_\text{pp}$] & $V_{A2}$ [V$_\text{pp}$] & $V_\text{OUT}$ [V$_\text{pp}$] & $V_\text{OUT}/V_S$ \\ \midrule
SK1 (bình thường) & 1,03 & 0,093 & 1,41 & 1,27 & 1236 \\
SK2 (nhanh, gần) & 1,85 & 0,170 & 2,45 & 2,23 & 1203 \\
SK3 (chậm, xa) & 0,49 & 0,045 & 0,61 & 0,54 & 1115 \\
\bottomrule
\end{tabularx}
\end{table}

Hệ số khuếch đại hiệu dụng giữa biên độ ra và biên độ vào là 1100 -- 1240 lần. Giá trị này nhỏ hơn một chút so với dải giữa (1413), vì tín hiệu chuyển động không phải hình sin: nó có thành phần tần số thấp gần tần số cắt 0,5~Hz. Kịch bản SK3 chậm nhất nên bị suy giảm nhiều nhất. Đầu ra tầng 1 chỉ dao động quanh 0,81~V trong khoảng 0,73 -- 0,90~V, không bị bão hòa, khẳng định việc chọn $K_1 = 101$.

\subsection{Toàn bộ các kịch bản}

Hình~\ref{fig:toanbo} là kết quả của cả năm kịch bản trong 40~s. Bảng~\ref{tab:kichban-kq} tổng hợp giá trị lớn nhất và nhỏ nhất của đầu ra.

\begin{figure}[!htb]
\centering
\begin{tikzpicture}
\begin{groupplot}[group style={group size=1 by 3,vertical sep=0.25cm,xticklabels at=edge bottom,xlabels at=edge bottom},
  width=0.93\textwidth,xmin=0,xmax=40,xtick distance=4,xlabel={Thời gian [s]},
  legend style={at={(0.0,1.0)},anchor=north west,legend columns=2}]
  \nextgroupplot[height=3.6cm,ylabel={$P$ [µW]},ymin=-0.1,ymax=2.3]
    \addplot[c1,thick] table[x=t,y=pa] {data/tran_full.dat}; \addlegendentry{$P_A$}
    \addplot[c2,thick] table[x=t,y=pb] {data/tran_full.dat}; \addlegendentry{$P_B$}
  \nextgroupplot[height=3.6cm,ylabel={$\Delta V_S$ [mV]},ymin=-1.2,ymax=1.4]
    \addplot[c3,thick] table[x=t,y=s] {data/tran_full.dat};
  \nextgroupplot[height=5.2cm,ylabel={$V_\text{OUT}$ [V]},ymin=-0.1,ymax=3.3,ytick={0,0.5,...,3}]
    \addplot[c1,thick] table[x=t,y=out] {data/tran_full.dat};
    \draw[bkred,densely dashed] (axis cs:0,3) -- (axis cs:40,3) node[pos=0.93,below,font=\scriptsize]{3,0 V};
    \draw[bkred,densely dashed] (axis cs:0,0) -- (axis cs:40,0);
    \draw[cgray,densely dashed] (axis cs:0,1.5) -- (axis cs:40,1.5);
    \node[font=\scriptsize] at (axis cs:4.8,2.85) {SK1};
    \node[font=\scriptsize] at (axis cs:11,2.85) {SK2};
    \node[font=\scriptsize] at (axis cs:16.8,2.85) {SK3};
    \node[font=\scriptsize] at (axis cs:24,2.85) {SK4};
    \node[font=\scriptsize] at (axis cs:35,2.85) {SK5};
\end{groupplot}
\end{tikzpicture}
\caption{Kết quả mô phỏng 40~s: tín hiệu vào, điện áp chân S và điện áp ra}
\label{fig:toanbo}
\end{figure}

\begin{table}[!htb]
\centering
\caption{Điện áp ra trong từng kịch bản}
\label{tab:kichban-kq}
\small
\begin{tabularx}{\textwidth}{@{}p{3.6cm}YYYL@{}}
\toprule
\textbf{Kịch bản} & \textbf{$V_\text{OUT,min}$ [V]} & \textbf{$V_\text{OUT,max}$ [V]} & \textbf{Đỉnh -- đỉnh [V]} & \textbf{Nhận xét} \\ \midrule
Nghỉ (0,5 -- 2,5~s) & \multicolumn{2}{c}{1,4992 (trung bình)} & 0,004 & Ổn định ở 1,5~V \\
SK1 bình thường & 0,80 & 2,07 & 1,27 & Dao động rõ \\
SK2 nhanh, gần & 0,32 & 2,55 & 2,23 & Lớn nhất, vẫn trong 0 -- 3~V \\
SK3 chậm, xa & 1,21 & 1,76 & 0,54 & Nhỏ nhưng vẫn phân biệt được \\
SK4 nền đổi đều & 1,46 & 1,53 & 0,07 & Gần như không đổi \\
SK5 vật đứng yên & 1,27 & 1,93 & 0,67 & Một xung rồi trở về 1,5~V \\
\bottomrule
\end{tabularx}
\end{table}

Trong mọi kịch bản, đầu ra luôn nằm trong khoảng \textbf{0,32 -- 2,55~V}, tức là trong dải 0 -- 3~V yêu cầu, và không bị bão hòa. Biên độ ra thay đổi theo cường độ chuyển động: người càng gần hoặc đi càng nhanh thì dao động càng lớn.

\subsection{Nguyên lý phần tử kép và vật đứng yên}
\label{sec:kq-sk45}

Hình~\ref{fig:sk45} phóng to hai kịch bản kiểm tra nguyên lý của cảm biến.

\begin{figure}[!htb]
\centering
\begin{tikzpicture}
\begin{groupplot}[group style={group size=2 by 3,horizontal sep=1.6cm,vertical sep=0.25cm,
    xticklabels at=edge bottom,xlabels at=edge bottom},
  width=0.5\textwidth,height=3.8cm,xlabel={Thời gian [s]}]
  \nextgroupplot[xmin=20.5,xmax=30.5,ylabel={$P$ [µW]},ymin=-0.1,ymax=2.2,title={\small SK4: nền tăng đều 2 µW}]
    \addplot[c1,thick] table[x=t,y=pa] {data/tran_sk4.dat};
    \addplot[c2,thick,dashed] table[x=t,y=pb] {data/tran_sk4.dat};
  \nextgroupplot[xmin=30,xmax=40,ymin=-0.01,ymax=0.14,title={\small SK5: vật nóng đứng yên}]
    \addplot[c1,thick] table[x=t,y=pa] {data/tran_sk5.dat};
    \addplot[c2,thick,dashed] table[x=t,y=pb] {data/tran_sk5.dat};
  \nextgroupplot[xmin=20.5,xmax=30.5,ylabel={$\Delta V_S$ [mV]},ymin=-0.4,ymax=1.3]
    \addplot[c3,thick] table[x=t,y=s] {data/tran_sk4.dat};
  \nextgroupplot[xmin=30,xmax=40,ymin=-0.4,ymax=1.3]
    \addplot[c3,thick] table[x=t,y=s] {data/tran_sk5.dat};
  \nextgroupplot[xmin=20.5,xmax=30.5,ylabel={$V_\text{OUT}$ [V]},ymin=1.0,ymax=2.0]
    \addplot[c1,thick] table[x=t,y=out] {data/tran_sk4.dat};
  \nextgroupplot[xmin=30,xmax=40,ymin=1.0,ymax=2.0]
    \addplot[c1,thick] table[x=t,y=out] {data/tran_sk5.dat};
\end{groupplot}
\end{tikzpicture}
\caption{Kiểm chứng nguyên lý: triệt tiêu thay đổi bức xạ chung (SK4, trái) và không phát hiện vật đứng yên (SK5, phải)}
\label{fig:sk45}
\end{figure}

\paragraph{SK4 -- triệt tiêu nhiễu chung.} Công suất tới tăng 2~µW trên \emph{cả hai} phần tử, lớn gấp 15 lần SK1. Nếu chỉ có một phần tử, theo \eqref{eq:step} điện áp chân S sẽ thay đổi khoảng 20~mV và đầu ra bão hòa hoàn toàn. Nhờ cấu trúc phần tử kép, $V_S$ chỉ thay đổi 0,34~mV$_\text{pp}$ (phần dư này do hai phần tử lệch nhau 1\%). Đầu ra chỉ dao động $\pm 35$~mV quanh 1,5~V. Độ triệt tiêu đạt khoảng 35 -- 40~dB, đúng với thông số match 1\% của datasheet.

\paragraph{SK5 -- vật đứng yên.} Công suất trên phần tử A tăng 0,12~µW rồi giữ nguyên. Theo \eqref{eq:step}, với $K_S = \num{39000}$~V·s/W, $\tau_E = 2{,}5$~s, $\tau_T = 0{,}6$~s:
\begin{equation*}
  t^* = \frac{0{,}6\times2{,}5}{2{,}5-0{,}6}\ln\frac{2{,}5}{0{,}6} = 1{,}13~\text{s},\qquad
  V_{S,\max} = \frac{K_S\,\Delta\Phi}{\tau_E-\tau_T}\left(e^{-t^*/\tau_E}-e^{-t^*/\tau_T}\right) \approx 1{,}19~\text{mV}.
\end{equation*}
Mô phỏng cho $V_S$ tăng 1,21~mV, khớp với tính toán, rồi giảm dần về mức ban đầu. Đầu ra có một xung ($+0{,}43$~V rồi $-0{,}23$~V), sau đó trở về $1{,}505$~V, dù vật nóng vẫn còn. Đây chính là nhược điểm ``không phát hiện người đứng yên'' của cảm biến PIR, được giải thích bằng nguyên lý nhiệt điện ở mục~\ref{sec:dapungbac}.

\subsection{Nhận biết hướng chuyển động}
\label{sec:kq-huong}

Em chạy thêm một kịch bản SK6: người đi \emph{ngược chiều} với SK1, cùng tốc độ và khoảng cách. Khi đó ảnh đi qua phần tử B trước, rồi mới tới A. Hình~\ref{fig:huong} so sánh đầu ra của hai kịch bản, đặt cùng gốc thời gian.

\begin{figure}[!htb]
\centering
\begin{tikzpicture}
\begin{axis}[width=0.88\textwidth,height=5.2cm,xmin=-0.5,xmax=4.5,ymin=0.6,ymax=2.75,
  xlabel={Thời gian kể từ lúc bắt đầu kịch bản [s]},ylabel={$V_\text{OUT}$ [V]},legend pos=north east]
  \addplot[c1,very thick] table[x=t,y=out] {data/dir_sk1.dat}; \addlegendentry{SK1: đi từ A sang B}
  \addplot[c2,very thick,dashed] table[x=t,y=out] {data/dir_sk6.dat}; \addlegendentry{SK6: đi từ B sang A}
  \draw[cgray,densely dashed] (axis cs:-0.5,1.5) -- (axis cs:4.5,1.5);
\end{axis}
\end{tikzpicture}
\caption{Đầu ra khi người đi theo hai hướng ngược nhau}
\label{fig:huong}
\end{figure}

Hai tín hiệu có cùng biên độ (1,27 và 1,28~V$_\text{pp}$) nhưng \emph{ngược pha}: SK1 tăng lên trước, SK6 giảm xuống trước. Như vậy, nếu phân tích thêm thứ tự xung dương/âm, hệ thống có thể biết được hướng đi của người.

\subsection{Hiệu quả của bộ lọc nhiễu}
\label{sec:kq-nhieu}

Hình~\ref{fig:nhieu} phóng to khoảng 0,5 -- 1,0~s, khi không có chuyển động. Ba đồ thị lần lượt là nhiễu đưa vào đường tín hiệu, tín hiệu trước bộ lọc Sallen-Key ($\Delta V_{A2} = V_{A2}-\VREF$) và sau bộ lọc ($\Delta V_{A3} = V_{A3}-\VREF$). Bảng~\ref{tab:nhieu} tổng hợp biên độ nhiễu đỉnh -- đỉnh đo trong 0,5 -- 2,5~s.

\begin{figure}[!htb]
\centering
\begin{tikzpicture}
\begin{groupplot}[group style={group size=1 by 3,vertical sep=0.25cm,xticklabels at=edge bottom,xlabels at=edge bottom},
  width=0.93\textwidth,height=3.5cm,xmin=0.5,xmax=1.0,xlabel={Thời gian [s]},
  legend style={at={(1.0,1.0)},anchor=north east}]
  \nextgroupplot[ylabel={[µV]},ymin=-130,ymax=130]
    \addplot[cgray,thin] table[x=t,y=n] {data/noise.dat}; \addlegendentry{$v_n$: nhiễu vào (50 Hz + 1 kHz + trắng)}
  \nextgroupplot[ylabel={[mV]},ymin=-12,ymax=12]
    \addplot[c2,thin] table[x=t,y=a2] {data/noise.dat}; \addlegendentry{$\Delta V_{A2}$: trước bộ lọc}
  \nextgroupplot[ylabel={[mV]},ymin=-12,ymax=12]
    \addplot[c1,thick] table[x=t,y=a3] {data/noise.dat}; \addlegendentry{$\Delta V_{A3}$: sau bộ lọc}
\end{groupplot}
\end{tikzpicture}
\caption{Nhiễu trước và sau bộ lọc thông thấp Sallen-Key khi không có chuyển động}
\label{fig:nhieu}
\end{figure}

\begin{table}[!htb]
\centering
\caption{Biên độ nhiễu khi không có chuyển động}
\label{tab:nhieu}
\small
\begin{tabularx}{\textwidth}{@{}p{4.2cm}p{3.0cm}L@{}}
\toprule
\textbf{Vị trí} & \textbf{Nhiễu đỉnh -- đỉnh} & \textbf{Thành phần chính} \\ \midrule
Đưa vào (SN $-$ S) & 0,22~mV & 50~Hz, 1~kHz, nhiễu trắng \\
Trước bộ lọc ($V_{A2}$) & 19,1~mV & 50~Hz ($\approx 15$~mV$_\text{pp}$) \\
Sau bộ lọc ($V_{A3}$) & 4,2~mV & Nhiễu trắng trong dải 0,5 -- 10~Hz \\
Đầu ra OUT & 3,8~mV & Bằng 0,13\% dải 3~V \\
\bottomrule
\end{tabularx}
\end{table}

Phân tích kết quả:
\begin{itemize}
  \item Trước bộ lọc, nhiễu 50~Hz được khuếch đại $10^{(29{,}7+13{,}7)/20} \approx 148$ lần, thành khoảng 15~mV$_\text{pp}$. Nó hiện rõ thành một dao động hình sin trên $V_{A2}$.
  \item Bộ lọc Sallen-Key làm suy giảm 50~Hz thêm 26,8~dB (khoảng 22 lần), nên dao động 50~Hz gần như biến mất trên $V_{A3}$. Nhiễu 1~kHz không còn nhận ra được.
  \item Phần nhiễu còn lại (khoảng 4~mV$_\text{pp}$) là \emph{nhiễu trắng nằm trong dải 0,5 -- 10~Hz}. Nó trùng tần số với tín hiệu nên không bộ lọc nào loại được. Mức nhiễu này nhỏ hơn nhiều so với tín hiệu nhỏ nhất (SK3: 0,54~V$_\text{pp}$).
  \item Cần lưu ý: chính cảm biến cũng có nhiễu riêng khoảng 20~µV$_\text{pp}$ trong dải 0,4 -- 10~Hz (datasheet). Sau khuếch đại, nhiễu này thành khoảng 25~mV$_\text{pp}$ ở đầu ra. Đây là giới hạn cơ bản quyết định ngưỡng phát hiện nhỏ nhất của hệ thống.
\end{itemize}

\subsection{Đánh giá kết quả so với yêu cầu}

\begin{table}[!htb]
\centering
\caption{Đánh giá kết quả mô phỏng so với yêu cầu thiết kế}
\label{tab:danhgia}
\small
\begin{tabularx}{\textwidth}{@{}p{4.3cm}p{3.4cm}Lc@{}}
\toprule
\textbf{Yêu cầu} & \textbf{Mục tiêu} & \textbf{Kết quả mô phỏng} & \textbf{Đạt} \\ \midrule
Mô phỏng đầu thu \LHI{} và các chân ra & Đúng nguyên lý, khớp datasheet & Độ nhạy 4,06~kV/W @1~Hz, $V_S = 0{,}81$~V, $r_\text{out} = 5{,}1$~kΩ & Đạt \\
Đáp ứng đúng nguyên lý cảm biến & Chỉ đáp ứng với thay đổi, vi sai & Nền đổi đều: $\pm35$~mV; vật đứng yên: trở về 1,5~V & Đạt \\
Nguồn cấp & 3,3~V đơn & Toàn mạch dùng nguồn 3,3~V & Đạt \\
Khuếch đại dùng op-amp & $\approx 60$~dB & 62,4~dB @1~Hz (AD8608) & Đạt \\
Lọc thông thấp dùng op-amp & Loại nhiễu tần số cao & Sallen-Key bậc 2; 50~Hz: $-47$~dB; 1~kHz: $-137$~dB & Đạt \\
Mức ra khi không có chuyển động & 1,5~V & 1,4992~V & Đạt \\
Dải ra khi có chuyển động & 0 -- 3~V & 0,32 -- 2,55~V; giới hạn cứng 3,0~V & Đạt \\
Một project LTspice duy nhất & 1 file & \texttt{PIR\_LHI778.asc} & Đạt \\
\bottomrule
\end{tabularx}
\end{table}
